% MAIN SNAPSHOT — LSPV/RMV::DIAGTRANSFER1::2026-10-13T13:52:14+09:00
% !TEX program = xelatex
\documentclass[11pt]{article}
\usepackage[T1]{fontenc}
\usepackage{lmodern,amsmath,amssymb,amsthm,booktabs,tabularx,xurl,hyperref}
\usepackage[margin=0.9in]{geometry}
\newtheorem{theorem}{Theorem}
\newtheorem{definition}[theorem]{Definition}
\newcommand{\Ltwo}{L_{22}}
\newcommand{\sCP}{\mathsf{sCP}}
\newcommand{\EXACT}{\textbf{PROVED/EXACT}}
\newcommand{\KNOWN}{\textbf{KNOWN THEOREM}}
\newcommand{\COMP}{\textbf{COMPUTATIONAL}}
\newcommand{\COND}{\textbf{CONDITIONAL}}
\newcommand{\CAND}{\textbf{CANDIDATE}}
\newcommand{\OPEN}{\textbf{OPEN}}
\newcommand{\RET}{\textbf{FALSIFIED/RETIRED}}
\title{DIAGTRANSFER1: Exponent-Indexed Families Are Sufficient;\\the Computation-to-Proof Arrow Is Not}
\author{UNIVERSAL\_CORE\_SYNTHESIS / LSPV-5}
\date{13 October 2026, 13:52 KST}

\begin{document}
\maketitle

\begin{center}
\fbox{\parbox{0.93\linewidth}{
\textbf{Executive status: P versus NP remains OPEN.}
The fixed target is satisfiability of valid nonempty proper distinct-clause
SIGNED exact-$(2,2)$ 3-CNFs under the declared binary encoding; malformed
and out-of-domain strings are rejected. ROUNDINGSQUEEZE5 is retained only as
a \CAND{} proof-model theorem pending independent mathematical review.

This run closes one quantifier objection exactly. A single family working for
all polynomial exponents is unnecessary: after fixing an alleged polynomial
decider and its exponent, one may choose a larger depth exponent and the
corresponding ROUNDINGSQUEEZE5 family. The composition still fails because no
verified theorem maps that decider to a bounded-size semantic-Cutting-Planes
proof or triangle DAG with polynomially related depth. Proof search,
verification, and one-bit decision remain different resources.}}
\end{center}

\section{Fixed language and exact target}
Let $\mathcal C_{22}$ be the valid nonempty proper distinct-clause signed
exact-$(2,2)$ 3-CNF domain, and
\[
 \Ltwo=\{\operatorname{enc}(G):G\in\mathcal C_{22},\ G\text{ satisfiable}\}.
\]
The project target is
\[
 \forall A\ \forall K\in\mathbb N\ \forall C>0\ \forall N_0\
 \exists x\quad
 [|x|\ge N_0\land x\text{ valid}\land
 T_A(x)>C|x|^K],                                         \tag{1}
\]
where $A$ ranges over correct total deterministic uniform deciders. Hard
inputs may depend on $A,K,C,N_0$. A fixed language is not a fixed formula,
and a family of certified NO instances is not by itself the decision
language.

The verified IMAGEEQUIV1 normal form permits replacing $\Ltwo$ by its total
recognizable ordered compiler-image language without changing membership in
$\mathrm P$. This is an algorithmic simplification, not a lower bound.

\section{Primary-literature preflight}
The current arXiv text of Mossel--Niles-Weed--Sun--Zadik states its normalized
spread lemma for an $R$-spread probability measure on sets of size at most
$k$, with hitting probability at least $0.9$ once
$p\ge C\log k/R$ \cite[Definition 1.1 and Theorem 1.2]{spread}. This is a
direct theorem about random set systems. It does not mention proof search or
arbitrary SAT machines.

The de Rezende--Fleming--Janett--Nordstr\"om--Pang preprint proves
supercritical trade-offs in named models: Resolution, semantic Cutting
Planes, monotone real circuits, and Weisfeiler--Leman \cite{tradeoff}. Its
Theorem 2.13 explicitly converts a semantic Cutting Planes refutation to a
triangle DAG before applying the lifting theorem. That direction is a
direct model-matched theorem for a proof already in semantic Cutting Planes;
it is not a conversion of an arbitrary deterministic decision computation
into such a proof.

Cook--Reckhow gives the exact framework for polynomially checkable
propositional proof systems and their simulations \cite{cr}. A decider can
define a tailored proof system, but there is no general simulation of every
such system by Resolution, PC/PCR, or semantic Cutting Planes. Treating that
simulation as automatic would insert the missing separation-strength claim.

\section{The proof-model input from ROUNDINGSQUEEZE5}
Write $\mathsf{RS5}(p,\beta)$ for the following reported statement. For fixed
rational $p>1$ and $1<\beta<1+1/p$, there are unsatisfiable
$H^{p,\beta}_t\in\mathcal C_{22}$, encoded lengths $L_t\to\infty$, and
polynomial budgets $r_t$ such that:
\begin{enumerate}
\item $H^{p,\beta}_t$ has a Resolution refutation with at most $r_t$ lines;
\item every semantic Cutting Planes refutation $\tau$ with
  $S(\tau)\le r_t^\beta$ has
  $D(\tau)\ge c_{p,\beta}L_t^p$ for all sufficiently large $t$.
\end{enumerate}
ROUNDINGSQUEEZE5 supplies a proposed proof through a normalized-spread
lifting theorem, no-weakening subclause extraction, and an explicit
selector-compiler Resolution lift. Its own status is \CAND{} pending
independent review. In particular, this synthesis does not promote it.

\section{New exact conditional composition}
\begin{definition}[Charged decider-to-$\sCP$ transfer]
A transfer has exponent map $q:\mathbb N\to\mathbb R_{>0}$ if, for every
total deterministic $\Ltwo$ decider $A$ running in time $O(N^K)$ and every
fixed admissible $(p,\beta)$, it externally constructs, for all sufficiently
large $t$, an $\sCP$ refutation $\tau_{A,p,\beta,t}$ of
$H^{p,\beta}_t$ satisfying
\[
 S(\tau_{A,p,\beta,t})\le r_t^\beta,
 \qquad
 D(\tau_{A,p,\beta,t})\le a_{A,p,\beta}L_t^{q(K)}.       \tag{2}
\]
The finite exponent $q(K)$ is independent of the later choice of $p$.
Acquisition of the source/compiler map, queries to $A$, replay,
serialization, coefficients, proof verification, and joint workspace are
charged. The construction may inspect $A$ as code, but it receives no hidden
advice depending on $t$.
\end{definition}

\begin{theorem}[Exponent-indexed diagonal composition]
Assume $\mathsf{RS5}(p,\beta)$ for every fixed admissible pair. If a charged
decider-to-$\sCP$ transfer as in (2) exists, then $\Ltwo\notin\mathrm P$.
The hard family may depend on the alleged decider's exponent; no single
family working for all exponents is required.
\end{theorem}
\begin{proof}
Suppose $A$ decides $\Ltwo$ in time $O(N^K)$. Choose any rational
$p>\max\{1,q(K)\}$ and any rational
$1<\beta<1+1/p$. Apply the transfer to the family supplied by
$\mathsf{RS5}(p,\beta)$. Equations (2) and the lower bound give, for all
large $t$,
\[
 c_{p,\beta}L_t^p
 \le D(\tau_{A,p,\beta,t})
 \le a_{A,p,\beta}L_t^{q(K)},
\]
which is impossible because $L_t\to\infty$ and $p>q(K)$. Thus no polynomial
decider exists. The family selection occurs after $A$ and $K$ are fixed,
matching the $\forall A\,\exists x_A$ order in (1).
\end{proof}

The theorem is \EXACT{} as a conditional composition lemma. It is not a
separation theorem because both its proof-model input and especially its
transfer premise remain unproved. The transfer is the sole genuinely new
load-bearing arrow. A premise saying only that ``every fast decider has a
short proof in some decider-specific Cook--Reckhow system'' is insufficient;
the lower bound applies to semantic Cutting Planes/triangle DAGs, not to an
arbitrary tailored system.

\section{Candidate attacked and smallest repair}
The candidate attacked was:
\[
 \text{``ROUNDINGSQUEEZE5 fails to address arbitrary polynomial exponents
 because its family depends on }p\text{.''}               \tag{3}
\]
Statement (3) is \RET{}. Family dependence on $p$ is compatible with (1).
The smallest correct repair is not to demand one universal family; it is to
demand the model-matched transfer (2), including the size budget.

The following stronger shortcuts remain invalid:
\begin{itemize}
\item \textbf{Trace-to-proof shortcut.} A one-bit rejecting computation is
  not syntactically an $\sCP$ refutation. Verifying a supplied proof and
  finding a decision answer are different tasks.
\item \textbf{Existential-proof shortcut.} Existence of a short Resolution
  upper proof does not imply that $A$ discovers it, or that its trace can be
  promoted to $\sCP$ within $r_t^\beta$.
\item \textbf{Wrapper shortcut.} DEAD, COPY, DELAY, INTERP, and RECOMP
  wrappers can erase any claim that the original run must visibly expose a
  selected compatibility object. An external extractor is allowed only with
  its full new execution charged.
\item \textbf{Representation shortcut.} Compact arithmetic circuits,
  Gaussian elimination, integer counting/flow, decoder-based preprocessing,
  bounded-treewidth dynamic programming, and alternative encodings need not
  expand into the monomials or inequalities used by a restricted proof.
\end{itemize}

\section{Bounded reproducible arithmetic audit}
The exact-rational script
\path{experiments/review/exponent_indexed_diagonal_audit.py}
(commit \texttt{4d606b4b0e5e41faad9388bdc297014e80e9127f}, blob
\texttt{7b48579e8cd48c9c1e53421c4960358729378b7c}, SHA-256
\texttt{8f319c3e99bfbecbd1cff755fe67a298bd4e778bea7bbe3ea8c6e65f3f098011})
checked surrogate transfer exponents $K=1,\ldots,64$ using
\[
 p=K+1,\qquad \beta=1+\frac1{2p}.
\]
It verified $p>K$, $p>1$, $1<\beta<1+1/p$, and exact polynomial dominance
for constant ratios $1,10,10^6$. It rejected the hostile boundaries $p=K$,
$\beta=1$, $\beta=1+1/p$, and a fixed $p=2$ for $K\ge2$. All 703 checks
passed in 0.008 seconds. The immutable summary is
\path{results/raw/review/2026-10-13_1352_exponent_indexed_diagonal_audit.json}
(commit \texttt{9ee0b70d1bee9f666fb6c5deedc969593d54c7e5}, blob
\texttt{e4ed5ad3c572f93a08848f1379eb79e9fbbc15bd}, SHA-256
\texttt{3fbd1f7ca72ceee0550f4f05d2354ac2be17558655800325e6c6f8d508349a8d}).
This evidence is \COMP{} and checks only the exponent arithmetic. It does
not test or support the missing computation-to-proof transfer.

\section{Current proof spine}
\begin{center}
\begin{tabularx}{\linewidth}{p{0.30\linewidth}p{0.17\linewidth}X}
\toprule
Arrow & Status & Current verified meaning\\
\midrule
Declared $\Ltwo$ interface and invalid-input rejection & \EXACT & One total
fixed decision language.\\
$\Ltwo$ NP membership/completeness and recognized image equivalence &
\EXACT & The image language is a complete normal form; source recovery is an
algorithmic advantage.\\
Compiled exact-$(2,2)$ Resolution-hard MU family & \EXACT/\KNOWN & Restricted
Resolution result with its graph/pairing hypotheses; historical variable
count converts to $2^{\Omega(L/\log L)}$.\\
Graph-colouring all-field expanded-monomial PC/PCR lower bound &
\KNOWN/\COND & Restricted algebraic proof endpoint; independent application
review remains pending.\\
ROUNDINGSQUEEZE5 $\forall(p,\beta)\exists H_t$ proof trade-off & \CAND &
Named proof-model statement pending independent mathematical review.\\
Exponent-indexed composition needs one family for every $p$ simultaneously &
\RET & The family may be selected after the alleged decider exponent.\\
Every polynomial decider yields (2) with a $p$-independent exponent map &
\OPEN & No trace-to-$\sCP$/triangle-DAG simulation theorem.\\
Whole-machine lower bound (1) & \OPEN & P versus NP remains open.\\
\bottomrule
\end{tabularx}
\end{center}

\section{Newest specialist reconciliation}
\paragraph{LSPV-1 / IMAGEEQUIV1.}
\EXACT{} (conditional on the previously stated ordered decoder theorem):
$\Ltwo\in\mathrm P$ iff the total recognizable compiler-image language is
in $\mathrm P$. This removes image recognition as an objection but supplies
no lower bound.

\paragraph{LSPV-2 / EQSPECTRAL1.}
The characteristic-zero spectral theorem gives an exact polynomial separator
for the connected balanced promise support. It retires that distributional
candidate, not arbitrary $\Ltwo$ inputs.

\paragraph{LSPV-3 / MODRANKFW1.}
The executable EQSPECTRAL1 classifier over the fixed prime $1{,}000{,}003$
has a proved false-positive family at $r=166667$ because the mode polynomial
factors as $(5z+1)(7z+35)$ modulo that prime. Exact rational/fraction-free
rank repairs the implementation. The characteristic-zero theorem survives;
the promise remains easy. This correction is adopted and the fixed-prime
all-$r$ implementation claim is \RET{}.

\paragraph{LSPV-4 / SOURCEARC-NOCONTRACT1.}
The natural decode, one-variable restrict, properize, balance, and recompile
route produces larger $\Ltwo$ children. It is \EXACT{} for that construction
and retires that ARC, not every possible direct argument.

\section{Barrier, SAME\_REALIZATION, and hostile controls}
\paragraph{SAME\_REALIZATION.}
The conditional theorem compares the same transferred $\sCP$ proof in both
inequalities. It does not assign a proof to the original run for free.
Any external construction is a new computation with acquisition, replay,
serialization, coefficient bit complexity, and workspace charged. RS5's
strip search is existential; its cost cannot be retroactively charged to
$A$.

\paragraph{Relativization and algebrization.}
The diagonal arithmetic is model-neutral and does not evade either barrier.
The missing transfer must be audited separately; a generic black-box claim
would be especially suspect in relativized worlds. No Natural-Proofs claim
is invoked because no constructive large circuit property is proved.

\paragraph{Easy-family and representation controls.}
Horn, 2SAT, affine/Tseitin Gaussian elimination, bounded treewidth,
independent products, padding, redundant clauses, alternative encodings,
and one-bit output remain admissible. A proposed transfer that maps these
easy computations to lower-bounded proof objects without explaining the
discrepancy is false or mis-scoped.

\paragraph{Post-quantum ledger.}
Schur squares, quadratic relations, syzygies, collision entropy, and binding
do not match the present computation-to-proof object: \textbf{FAILED
TRANSFER}. Worst-case/average-case architecture is an
\textbf{ANALOGY/CANDIDATE}; it would still require a construction-specific,
unconditional black-box theorem. The exact exponent-selection lemma is an
\textbf{EXACT OBJECT/THEOREM TRANSFER} only after a model-matched proof
family and transfer have been supplied.

\section{Corpus reconciliation record}
\paragraph{CORPUS REVIEWED.}
The actual current \path{main4.tex} ARC driver and \path{main5.tex} WRAPPER
driver, repository protocol and research skill, task/governance ledgers,
EARLY22-INTEGRATION1, SINGLE22-SEPARATION1, the malformed-root incident,
BALMIXFW1, ROUNDINGSQUEEZE5, the newest four specialist reports, their cited
experiments, and the relevant compiler/proof-transfer history were reviewed.
Primary literature was checked before using spread, lifting, semantic CP,
or Cook--Reckhow terminology.

\paragraph{DUPLICATION / CONTRADICTION AUDIT.}
ROUNDINGSQUEEZE5 itself notes $\forall(p,\beta)\exists H_t$ rather than one
family for all $p$, but does not state the present conditional diagonal
composition against a general transfer exponent map. The new lemma does not
contradict TRACECR-FIREWALL1; it makes that report's missing promotion arrow
the explicit premise. MODRANKFW1 corrects only the executable finite-field
classifier and not EQSPECTRAL1's exact theorem.

\paragraph{WHAT CHANGED RELATIVE TO ROOT MAIN.}
The root was a documented 891-byte markerless placeholder. As whole-document
owner, this run restores the exact last verified full body from commit
\texttt{79b0501ca31c76a1529cb67b901fa6d5f21cff53} and adds a delimited current
synthesis section. No scientific body is reconstructed from the placeholder.
The scientific change is that ``one family for every exponent'' is removed
from the open-obligation list; the charged computation-to-$\sCP$ transfer is
retained as OPEN.

\paragraph{ADOPTED / REJECTED / RETIRED PRIOR COMPONENTS.}
Adopted: fixed $\Ltwo$, image equivalence, exact compiler identities,
restricted proof-model results in their scopes, EQSPECTRAL1 over
characteristic zero, and the SOURCEARC negative theorem. Rejected: any
unproved arbitrary-trace promotion, finite-field rank as exact rank, and
runtime inference from proof depth. Retired: the fixed-prime all-$r$
classifier, the naive balanced spectral distribution, the natural
source-first one-variable ARC, and objection (3).

\paragraph{STALE OR OVERSTATED CLAIMS FOUND.}
Any statement that ROUNDINGSQUEEZE5 needs one universal family for all
polynomial exponents is stale. Any statement that its proof-model theorem
already yields (1), or that the fixed-prime executable equals exact rank for
all $r$, is overstated.

\paragraph{EARLY22-INTEGRATION1 acknowledgment.}
\textbf{Integrated.} The original ledger and its cited exact
compiler/minimality/projection bodies and supersession notices remain
separate from the monotone cubic-pair result and affine-continuation
limitation. Historical $N$ is a variable count, not encoded bits. Exact
locations here are the proof-spine table and this acknowledgment. Next
check: independently verify RS5's saved-strip chronology and same-row
quantifier without reopening the completed exact-$(2,2)$ projection.

\paragraph{SINGLE22-SEPARATION1 acknowledgment.}
\textbf{Integrated.} The exact fixed-language priority is used throughout.
The arrow advanced is
\[
 \bigl[\forall p\ \exists H^{p}\bigr]
 +\text{ charged exponent-bounded transfer}
 \Longrightarrow \Ltwo\notin\mathrm P,
\]
with no family-broadening requirement. The transfer is OPEN, so the
single-language separation remains OPEN.

\section{Highest-value next target}
Independently audit ROUNDINGSQUEEZE5's spread lifting. In parallel, formulate
the weakest non-tautological transfer for a \emph{specific} restricted class
of $\Ltwo$ deciders---for example, deciders whose rejecting computations
already carry locally checkable integer inequalities---and prove or falsify
that transfer with exact size, depth, coefficient, and workspace costs.
Do not call it universal. A successful restricted transfer would identify
which additional semantic feature, if any, could plausibly be generalized;
a failure would prevent another hidden trace-to-proof assumption.

\begin{thebibliography}{9}
\bibitem{spread}
E. Mossel, J. Niles-Weed, N. Sun, and I. Zadik,
``A Second Moment Proof of the Spread Lemma,'' arXiv:2209.11347,
Definition 1.1 and Theorem 1.2. \url{https://arxiv.org/abs/2209.11347}.
\bibitem{tradeoff}
S. F. de Rezende, N. Fleming, D. A. Janett, J. Nordstr\"om, and S. Pang,
``Truly Supercritical Trade-offs for Resolution, Cutting Planes, Monotone
Circuits, and Weisfeiler--Leman,'' arXiv:2411.14267v1, 2024,
Theorems 2.3, 2.11, and 2.13.
\url{https://arxiv.org/abs/2411.14267}.
\bibitem{cr}
S. A. Cook and R. A. Reckhow,
``The Relative Efficiency of Propositional Proof Systems,''
\emph{Journal of Symbolic Logic} 44(1):36--50, 1979.
\url{https://www.cs.toronto.edu/~sacook/homepage/cook_reckhow.pdf}.
\end{thebibliography}

\end{document}
% END_OF_VERIFIED_MAIN::DIAGTRANSFER1
